% 基于既有简介补齐共用资料；身份与联系字段统一来自 metadata.tex。
\subsubsection{基本情况与联系方式}
申请人\CiesmValue{applicant}，\CiesmValue{sex}，\CiesmValue{institution}\CiesmValue{professional-title}、\CiesmValue{degree}、\CiesmValue{mentor}，已开展多尺度时间序列、能源预测和多源信息融合研究。

通讯地址：\CiesmValue{address}。

联系电话：\CiesmValue{phone}。专业邮箱：\CiesmValue{email}。

英文姓名\CiesmValue{english-name}；院系为\CiesmValue{faculty}；人才类别为\CiesmValue{talent}；办公地点为\CiesmValue{office}；个人主页\CiesmUrlValue{homepage}。

\subsubsection{教育经历}

\CiesmValue{education}

\subsubsection{研究工作简历}

\CiesmValue{employment}

申请人创建新疆大学 ALPHA Group（Adaptive Long-horizon Prediction with Hybrid AI），长期开展金融与能源时序建模、新能源发电功率预测及多源时序信息融合研究。

\subsubsection{与本课题相关的主要论著}
与本课题直接相关的多尺度时序建模代表作如下，相关研究积累及其在本课题中的作用见工作基础栏目。

\paragraph{DWT-Former}
\CiesmFullCite{DWT2025}。本人为通讯作者。

\paragraph{WaveKAN}
\CiesmFullCite{WaveKAN2026}。本人为共同第一作者。

\subsubsection{本课题分工}

申请人在本课题中负责总体设计、目标场信息含量检验、TerraSR-XJ 约束投影模块实现、验证方法设计、进度管理与成果组织。参与研究的 7 名研究生均就读于新疆大学：\mbox{孔哲}、\mbox{张建宏}拟负责资料处理与质量控制；\mbox{杨铄}、\mbox{孟星宇}、\mbox{桑乐}拟负责模型实验与消融；\mbox{吕中原}、\mbox{于得武}拟负责集成评价与接口联调。研究生在申请人指导下协同完成各项任务，劳务补助从本课题 1.10 万元劳务费预算中安排，并按实际参与工作量分配。
